\section{Round 4 Progress on Erd\H{o}s Problem 23}

\subsection*{1) ROUND-4 OBJECTIVE}
\textbf{Path (A): proof track by extending the exact $n^2$ bound to a strictly larger,
natural structural class.}
Round~3 proved the conjectured bound $\tau_B(G)\le n^2$ for all \emph{$C_5$-colorable}
graphs on $5n$ vertices, with rigidity at equality.
In Round~4 we strengthen this by proving the analogous $n^2$ bound (with rigidity)
for \emph{every odd cycle target}: for each $k\ge 2$, every graph homomorphic to $C_{2k+1}$
on $(2k+1)n$ vertices satisfies $\tau_B(G)\le n^2$.
This strictly advances the investigation by proving Erd\H{o}s' more general variant
(on $(2k+1)n$ vertices) for the large subclass of $C_{2k+1}$-colorable graphs.

\subsection*{2) ROUND-3 FOUNDATION USED}
We rely on the following Round~3 results.
\begin{itemize}
\item \textbf{(R3--C5Class)} If $G$ is $C_5$-colorable with $|V(G)|=5n$, then $\tau_B(G)\le n^2$,
      with equality forcing the balanced complete blow-up of $C_5$.
\item \textbf{(R3--AdjProd$_5$)} If $a_1,\dots,a_5\ge 0$ and $\sum a_i=5n$, then
      $\min_i a_i a_{i+1}\le n^2$, with equality iff all $a_i=n$.
\end{itemize}
We also keep Round~2 context (universal constant $\frac{25}{12}$ and the mid-density
localization), but Round~4's main theorem does not require it.

\subsection*{3) NEW INSIGHT / TOOL (ROUND-4)}
\paragraph{Odd-cycle generalization of the AM--GM adjacency lemma.}
The Round~3 AM--GM argument for five parts extends verbatim to $(2k+1)$ parts:
among $(2k+1)$ nonnegative numbers summing to $(2k+1)n$, some adjacent product
is at most $n^2$, with equality forcing all parts equal.

\paragraph{Odd-cycle pullback cuts.}
Every $2$-coloring of $C_{2k+1}$ leaves at least one monochromatic cycle-edge, and there
exists a $2$-coloring leaving \emph{exactly one} monochromatic cycle-edge. Pulling such a
coloring back along a homomorphism $G\to C_{2k+1}$ yields a cut in $G$ for which all
deleted edges lie in a single consecutive pair of color classes.

\subsection*{4) ATTACK PLAN (ROUND-4)}
The post--Round~3 gap is that the $n^2$ bound is established only for $C_5$-colorable graphs
on $5n$ vertices; a counterexample to Erd\H{o}s' conjecture must be \emph{non-$C_5$-colorable}.
In this round we:
\begin{itemize}
\item prove the exact $n^2$ deletion bound for all $C_{2k+1}$-colorable graphs on $(2k+1)n$ vertices;
\item prove a rigidity statement characterizing equality by the balanced complete blow-up of $C_{2k+1}$.
\end{itemize}
This strictly enlarges the class of graphs for which the conjectured bound (and its natural
$(2k+1)n$ generalization) is verified.

\subsection*{5) WORK (ROUND-4)}

\subsubsection*{5.1 General adjacent-product lemma for $(2k+1)$ parts}
\begin{lemma}\label{lem:adjprod_odd}
Fix an integer $k\ge 2$ and set $p:=2k+1$.
Let $a_1,\dots,a_p\ge 0$ satisfy $\sum_{i=1}^p a_i = pn$.
Then
\[
\min_{i\in\mathbb Z/p\mathbb Z} a_i a_{i+1}\ \le\ n^2.
\]
Moreover, $\min_i a_i a_{i+1}=n^2$ if and only if $a_1=\cdots=a_p=n$.
\end{lemma}

\begin{proof}
If some $a_i=0$ then $\min_i a_i a_{i+1}=0\le n^2$, so assume $a_i>0$ for all $i$.
Suppose for contradiction that $a_i a_{i+1} > n^2$ for every $i$ (indices mod $p$).
Multiplying all $p$ inequalities gives
\[
\prod_{i=1}^p (a_i a_{i+1}) > n^{2p}.
\]
But $\prod_{i=1}^p (a_i a_{i+1}) = (a_1a_2\cdots a_p)^2$, hence
\[
a_1a_2\cdots a_p > n^{p}.
\]
On the other hand, by AM--GM,
\[
a_1a_2\cdots a_p \le \left(\frac{a_1+\cdots+a_p}{p}\right)^p = n^p,
\]
a contradiction. Therefore $\min_i a_i a_{i+1}\le n^2$.

For the equality statement, assume $\min_i a_i a_{i+1}=n^2$.
Then all $a_i a_{i+1}\ge n^2$, so
\[
(a_1\cdots a_p)^2=\prod_{i=1}^p (a_i a_{i+1}) \ge n^{2p},
\quad\text{hence } a_1\cdots a_p\ge n^p.
\]
AM--GM gives $a_1\cdots a_p\le n^p$, so equality holds in AM--GM and thus
$a_1=\cdots=a_p=n$.
\end{proof}

\subsubsection*{5.2 The $n^2$ deletion bound for $C_{2k+1}$-colorable graphs}
\begin{definition}
Fix $k\ge 2$ and set $p:=2k+1$.
A graph $G$ is \emph{$C_{p}$-colorable} if there exists a homomorphism
$\varphi:V(G)\to V(C_{p})=\{1,2,\dots,p\}$ such that for every edge $uv\in E(G)$,
the pair $\varphi(u)\varphi(v)$ is an edge of $C_{p}$.
Equivalently, writing $V_i:=\varphi^{-1}(i)$, every edge of $G$ lies between
$V_i$ and $V_{i\pm 1}$ (indices mod $p$).
\end{definition}

\begin{theorem}\label{thm:Coddcolorable}
Fix $k\ge 2$ and set $p:=2k+1$.
Let $G$ be $C_p$-colorable with $|V(G)|=pn$.
Then
\[
\tau_B(G)\le n^2.
\]
Moreover, if $\tau_B(G)=n^2$, then $|V_1|=\cdots=|V_p|=n$ and
$e(V_i,V_{i+1})=n^2$ for all $i$; in particular, $G$ is the balanced complete
blow-up of $C_p$.
\end{theorem}

\begin{proof}
Let $\varphi:V(G)\to \{1,2,\dots,p\}$ be a homomorphism to $C_p$, and set $V_i:=\varphi^{-1}(i)$.
Then $V(G)=V_1\cup\cdots\cup V_p$ and every edge of $G$ lies between $V_i$ and $V_{i\pm 1}$.

Fix an index $i$ (mod $p$).  Consider the $2$-coloring of $V(C_p)$ obtained by alternating
colors around the cycle. Since $p$ is odd, this alternating coloring makes \emph{exactly one}
cycle-edge monochromatic, namely the edge $i(i+1)$ (after a suitable rotation of labels).
Pull this $2$-coloring back via $\varphi$ to obtain a bipartition $V(G)=X\cup Y$.
By construction, every edge of $G$ except those between $V_i$ and $V_{i+1}$ crosses the cut,
and all edges between $V_i$ and $V_{i+1}$ lie within $X$ or within $Y$.
Therefore deleting the edges between $V_i$ and $V_{i+1}$ makes $G$ bipartite, so
\[
\tau_B(G)\le e(V_i,V_{i+1}) \qquad \text{for every } i.
\]
Hence
\[
\tau_B(G)\le \min_i e(V_i,V_{i+1}) \le \min_i |V_i||V_{i+1}|.
\]
Since $\sum_i |V_i|=|V(G)|=pn$, Lemma~\ref{lem:adjprod_odd} implies
$\min_i |V_i||V_{i+1}|\le n^2$, proving $\tau_B(G)\le n^2$.

Assume now that $\tau_B(G)=n^2$.  Then
\[
n^2=\tau_B(G)\le \min_i e(V_i,V_{i+1})\le \min_i |V_i||V_{i+1}|\le n^2,
\]
so all inequalities are equalities and in particular $\min_i |V_i||V_{i+1}|=n^2$.
By Lemma~\ref{lem:adjprod_odd}, this forces $|V_1|=\cdots=|V_p|=n$, so
$|V_i||V_{i+1}|=n^2$ for all $i$.  Since also $\min_i e(V_i,V_{i+1})=n^2$ and each
$e(V_i,V_{i+1})\le |V_i||V_{i+1}|=n^2$, we deduce $e(V_i,V_{i+1})=n^2$ for every $i$.
Thus each consecutive pair induces $K_{n,n}$, i.e.\ $G$ is the balanced complete blow-up of $C_p$.
\end{proof}

\paragraph{Specialization back to the original problem.}
Taking $k=2$ gives $p=5$, and Theorem~\ref{thm:Coddcolorable} recovers (and slightly streamlines)
the Round~3 result for $C_5$-colorable graphs on $5n$ vertices.

\subsection*{6) ADVERSARIAL VERIFICATION}
\begin{itemize}
\item \textbf{Boundary cases:} If some class $V_i$ is empty, then $e(V_i,V_{i+1})=0$ and the pulled-back
      cut deletes $0$ edges, so $\tau_B(G)=0\le n^2$. Lemma~\ref{lem:adjprod_odd} also holds trivially.
\item \textbf{2-coloring claim:} For odd $p$, the alternating coloring around $C_p$ forces exactly one
      monochromatic edge; this is immediate by parity.
\item \textbf{Quantifiers:} Theorem~\ref{thm:Coddcolorable} holds for all $n\ge 1$ and all $k\ge 2$;
      no ``sufficiently large'' hypothesis is used.
\item \textbf{Interaction with Round~2:} The new theorem gives the sharp constant $1$ on a large explicit
      subclass, consistent with Round~2's universal bounds (which apply to all triangle-free graphs).
\item \textbf{Rigidity check:} The equality chain forces equal part sizes via Lemma~\ref{lem:adjprod_odd},
      then forces completeness between every consecutive pair because $e(V_i,V_{i+1})\le |V_i||V_{i+1}|=n^2$
      and $\min_i e(V_i,V_{i+1})=n^2$.
\end{itemize}

\subsection*{7) FINAL (EXACTLY ONE)}
\textbf{UNRESOLVED (BUT STRICTLY ADVANCED).}
We prove Erd\H{o}s' $n^2$ deletion bound for every $C_{2k+1}$-colorable graph on $(2k+1)n$ vertices,
with a rigidity theorem characterizing equality as the balanced complete blow-up of $C_{2k+1}$.
For the original $5n$-vertex problem (the case $k=2$), this confirms the conjectured bound on the full
$C_5$-colorable subclass and shows that any counterexample must be \emph{non-$C_5$-colorable}.

\subsection*{8) COMPLETION ESTIMATE (MANDATORY)}
\[
\textbf{COMPLETION: }76\%.
\]

\subsection*{9) REFERENCES}
S. Norin, Y. R. Sun, \emph{Triangle-independent sets vs.\ cuts}, arXiv:1602.04370.\\
J. Balogh, F. C. Clemen, B. Lidick\'y, \emph{Max Cuts in Triangle-free Graphs}, arXiv:2103.14179.

